\title{FlashAttention-3: Fast and Accurate Attention with Asynchrony and Low-precision}

\begin{abstract}
Attention mechanisms represent the central computational component within the Transformer framework, frequently constituting the primary performance bottleneck for both large language models and tasks requiring extensive context lengths. \textsc{FlashAttention} introduced a method to accelerate attention computation on GPUs by substantially reducing the frequency of memory access operations. Nevertheless, this approach has not fully leveraged features introduced in more recent GPU hardware, as evidenced by \textsc{FlashAttention-2} achieving merely 35\% of peak capacity on H100 GPUs. To further optimize attention performance on Hopper GPUs, we propose three core strategies: harnessing the asynchronous execution of Tensor Cores and TMA to (1) concurrently orchestrate computation and data transfer through warp-specialization and (2) alternate between block-level matrix multiplication and softmax steps, as well as (3) integrating block-wise quantization and incoherent processing by employing hardware-based FP8 low-precision support. Our experimental results show that \textsc{FlashAttention-3} attains 1.5-2.0$\times$ higher throughput on H100 GPUs, reaching up to 740 TFLOPs/s (75\% of theoretical capacity) for FP16, and approaching 1.2 PFLOPs/s for FP8 computations. We further establish that FP8 \textsc{FlashAttention-3} delivers numerical errors that are 2.6$\times$ lower in comparison to a standard FP8 attention baseline.
\end{abstract}

\section{Introduction}
\label{sec:intro}

Within the Transformer architecture~\citep{vaswani2017attention}, the attention module represents the central limitation in computational efficiency, as the calculation of self-attention scores between queries and keys incurs a quadratic cost relative to sequence length. Enhancing the scalability of attention mechanisms to process longer sequences promises to enable new functionalities—such as reasoning across multiple extensive documents~\citep{guo2021longt5,shaham2022scrolls,peng2023yarn}, managing numerous files in vast code repositories~\citep{roziere2023code, li2023starcoder}, handling emergent modalities like high-resolution images~\citep{chen2022scaling}, audio~\citep{gulati2020conformer}, and video~\citep{ho2022video}, as well as supporting advanced applications that rely on extensive user histories~\citep{sun2019bert4rec} or agent operations spanning prolonged horizons~\citep{yao2022react}. This need has stimulated extensive research aimed at accelerating attention in the context of longer sequences, pursued through approximation methods~\citep{katharopoulos2020transformers,choromanski2020rethinking, tay2020efficient}, advancements in software engineering (\citep{rabe2021self,dao2022flashattention,kwon2023efficient}), and the development of novel architectures~\citep{peng2023rwkv,sun2023retentive,gu2023mamba}.

This paper extends the efforts of \citet{dao2022flashattention}, who pioneered the design of exact-attention methods that strategically incorporate the GPU's architectural and computational model into algorithm development. In their contribution, Dao et al. introduced \textsc{FlashAttention}, presenting a novel tiling-based parallelization scheme for the attention mechanism, which merges all associated operations into a single GPU kernel and removes redundant data transfers with slow global memory. 

In subsequent work, \citet{dao2023flashattention2} refined the approach with \textsc{FlashAttention-2}, which enhanced the original by enabling parallel processing across the sequence length and restructuring the forward-pass inner loop to operate over partitions of the key and value matrices. These improvements led to better distribution of workload and resource occupation on the GPU. Nevertheless, our analysis reveals that \textsc{FlashAttention-2} still fails to fully utilize the potential of recent GPU architectures, particularly in comparison to highly-tuned GEMM routines, with utilization numbers of approximately 35\% as opposed to the 80-90\% typical for the Hopper H100 GPU. This shortfall can be partially traced back to differences at the implementation level—for instance, not leveraging Hopper-specialized instructions for Tensor Cores, which instead default to instructions from the previous Ampere generation. Recent advances such as ThunkerKitten~\citep{spector2024thunder} and cuDNN 9~\citep{cudnn9} demonstrate that employing Hopper-tailored instructions and leveraging tiling-based abstraction can both accelerate attention computation and streamline the implementation process.

At its core, the algorithm underlying \textsc{FlashAttention-2} operates within a basic synchronous execution framework and does not directly leverage hardware-level asynchrony or low-precision arithmetic in its formulation. In modern hardware, asynchrony is inherent due to the presence of specialized units designed to accelerate dominant ML operations: Tensor Cores are devoted to matrix multiplications, while other subsystems, such as the Tensor Memory Accelerator (TMA), handle memory transfers independently of general CUDA cores responsible for standard logical, integer, or floating-point tasks. The adoption of lower-precision formats, including FP8 in Hopper and FP4 in Blackwell—following on from FP16 (introduced with Pascal in 2017) and BF16 (with Ampere in 2020)—has demonstrably increased computational throughput (by up to two or four times) within fixed energy and area budgets. We discuss these advancements enabled by Hopper’s architecture in greater detail in \cref{subsec:hardware}. The primary engineering hurdle lies in refactoring \textsc{FlashAttention-2} to take advantage of these architectural features: asynchrony implies enabling the concurrent execution of matmul and softmax computations, even as their data dependencies are nontrivial, while deploying low-precision arithmetic calls for meticulous strategies to control quantization errors, particularly for handling outlier elements prevalent in LLMs~\citep{dettmers2208llm, sun2024massive}.

Motivated by these considerations, we introduce \textsc{FlashAttention-3}, presenting a unified framework of three key innovations to advance performance for state-of-the-art GPU platforms.\footnote{While our implementation and experiments are set within the NVIDIA Hopper ecosystem, our proposed algorithm is applicable to any GPU platform offering mature support for asynchronous execution and low-precision data paths.}

\begin{enumerate}

\item \textbf{Producer-Consumer asynchrony:} We propose a warp-specialized approach to software pipelining that takes advantage of the asynchronous nature of data transfers and Tensor Core computations. By assigning data producers and consumers to different warps, our method broadens the scope for hiding both memory access and instruction issue latency within the algorithm.
\item \textbf{Hiding softmax under asynchronous block-wise GEMMs:} To mask the performance overhead from the low-throughput, non-GEMM softmax components (like floating-point multiply-add and exponentiation), we interleave their execution with the asynchronous WGMMA-based GEMM steps. In redesigning the \textsc{FlashAttention-2} algorithm, we eliminate certain stepwise dependencies between softmax operations and GEMMs. Specifically, in the two-stage version, while softmax operates over one scores matrix block, WGMMA concurrently computes the subsequent block through asynchronous scheduling.
\item \textbf{Hardware-accelerated low-precision GEMM:} Our modifications to the forward pass enable execution on FP8 Tensor Cores, delivering nearly a twofold increase in observed TFLOPs/s. This adaptation necessitates resolving mismatches between WGMMA's required memory layouts for FP8 operands and FP32 accumulators. We employ block quantization techniques and support incoherent computation to address and limit precision degradation that arises from the use of FP8 arithmetic.
\end{enumerate}

To empirically assess the effectiveness of our approach, we evaluate \textsc{FlashAttention-3} on the H100 SXM5 GPU across various parameter settings. Our results demonstrate that (1) FP16 mode delivers a 1.5–2.0$\times$ acceleration over \textsc{FlashAttention-2} in the forward computation (peaking at 740 TFLOPs/s), and achieves a 1.5–1.75$\times$ improvement in the backward computation, (2) FP8 mode nearly reaches 1.2 PFLOPs/s, and (3) for tasks involving long sequences, FP16 consistently leads and FP8 attains performance on par with, or even surpassing in specific cases,\footnote{More precisely, for head dimension 64 \textsc{FlashAttention-3} FP8 is ahead, while for head dimensions 128 and 256 it is at par for those cases without causal masking and behind with causal masking.} the state-of-the-art attention implementation provided by NVIDIA’s cuDNN library. Further analysis indicates that FP16 \textsc{FlashAttention-3} matches \textsc{FlashAttention-2} in terms of numerical precision, and exceeds standard attention routines where intermediate values (such as softmax scaling) utilize FP32. Additionally, FP8 \textsc{FlashAttention-3} employing block quantization and incoherent processing is 2.6$\times$ as accurate as standard attention using per-tensor quantization, particularly in scenarios containing outlier features.

We are releasing \textsc{FlashAttention-3} under a permissive open-source license\footnote{\textsc{FlashAttention-3}
  is available at \texttt{https://github.com/Dao-AILab/flash-attention}} and are committed to integrating this implementation with both the PyTorch and Hugging Face ecosystems to maximize accessibility for researchers and practitioners.

\section{Background: Multi-Head Attention and GPU Characteristics}
\label{sec:background}

\subsection{Multi-Head Attention}
\label{subsec:multi_head_attn}

Consider the query, key, and value input sequences $\mathbf{Q}, \mathbf{K}, \mathbf{V} \in \mathbb{R}^{N \times d}$ for a single attention head, where $N$ represents the number of tokens in the sequence and $d$ denotes the dimensionality of each head. The computation of the attention output $\mathbf{O}$ proceeds as follows:

\begin{equation*}
  \mathbf{S} = \alpha \mathbf{Q} \mathbf{K}^\top \in \mathbb{R}^{N \times N}, \quad \mathbf{P} = \mathrm{softmax}(\mathbf{S}) \in \mathbb{R}^{N \times N}, \quad \mathbf{O} = \mathbf{P}\mathbf{V} \in \mathbb{R}^{N \times d},
\end{equation*}

Here, the $\mathrm{softmax}$ function operates across each row, and it is common to choose $\alpha = 1/\sqrt{d}$ as the scale parameter. To mitigate issues related to numerical instability during exponentiation, we subtract $\mathrm{rowmax}(\mathbf{S})$ from $\mathbf{S}$ prior to computing the softmax. In the context of multi-head attention (MHA), separate projections for queries, keys, and values are maintained for each head, with these computations executed in parallel over all heads and batches, yielding the complete output tensor.

Consider $\phi$ as a scalar-valued loss function and let us define $\mathbf{d}(-) = \partial \phi / \partial (-)$ to represent the gradient operation. Provided the gradient of the outputs, $\mathbf{dO} \in \mathbb{R}^{N \times d}$, the gradients with respect to $\mathbf{Q}$, $\mathbf{K}$, and $\mathbf{V}$ can be obtained via the chain rule as shown:

\begin{align*}
  \mathbf{dV} &= \mathbf{P}^\top \mathbf{dO} \in \mathbb{R}^{N \times d} \\
  \mathbf{dP} &= \mathbf{dO} \mathbf{V}^\top \in \mathbb{R}^{N \times N} \\
  \mathbf{dS} &= \mathrm{dsoftmax} (\mathbf{dP}) \in \mathbb{R}^{N \times N} \\
  \mathbf{dQ} &= \alpha \mathbf{dS} \mathbf{K} \in \mathbb{R}^{N \times d} \\
  \mathbf{dK} &= \alpha \mathbf{dS}^\top \mathbf{Q} \in \mathbb{R}^{N \times d},
\end{align*}

For the $\mathrm{softmax}$ function applied to a vector $s$, where $p = \mathrm{softmax}(s)$, we have $\mathbf{d}s = (\mathrm{diag}(p) - p p^\top)\mathbf{d}p$. The notation $\mathrm{dsoftmax}(\mathbf{dP})$ denotes that this transformation is executed independently for each row. Importantly, this approach can be computed in parallel for each head and batch during the backward computation step in multi-head attention (MHA).

\subsection{GPU hardware characteristics and execution model}
\label{subsec:hardware}

We describe the aspects of the GPU's execution model relevant for \textsc{FlashAttention-3}, with a focus on the NVIDIA Hopper architecture as a concrete instantiation of this model.

\paragraph{Memory hierarchy:} The GPU features a layered structure of memory locations, where higher bandwidth corresponds to smaller storage capacities (\cref{tab:gpu-hierarchy})\footnote{\citet{luo2024benchmarking} indicates that the shared memory bandwidth reaches 128 bytes per clock cycle for each SM; multiplying by 132 SMs and a boost clock speed of 1830 MHz yields the aggregate value.}. Global memory (GMEM)—also referred to as HBM—serves as the off-chip DRAM that can be accessed by all streaming multiprocessors (SMs). Data stored in GMEM can be automatically cached in the on-chip L2 cache. Beyond this, every SM integrates its own on-chip shared memory (SMEM), which is a small, highly banked cache under the explicit control of programmers. The lowest level in this hierarchy is represented by the register file available within each SM.

\paragraph{Thread hierarchy:} In the GPU programming framework, execution units are systematically arranged into hierarchical groups referred to as threads. This hierarchy progresses from the smallest units—individual threads—to increasingly larger groupings: warps (collections of 32 threads), warpgroups (consisting of 4 consecutive warps), threadblocks (also known as cooperative thread arrays or CTAs), threadblock clusters (introduced in Hopper), and, at the broadest scale, grids.

The two hierarchies are tightly connected. Threads belonging to a single CTA are executed together on a single SM, and CTAs assigned to one cluster are scheduled jointly on the same GPC. All threads within a CTA have direct access to the SMEM, while each thread can utilize up to 256 dedicated registers (RMEM) exclusively.

\paragraph{Asynchrony and warp-specialization:}

GPUs function as throughput-oriented processors, leveraging both concurrent execution and asynchronous operations to mask latencies in memory access and instruction execution. In the case of asynchronous data transfer between GMEM and SMEM, Hopper introduces the Tensor Memory Accelerator (TMA), a specialized hardware component \cite[\S7.29]{cuda}. Additionally, diverging from earlier designs like Ampere, Hopper’s Tensor Core—accessed with the warpgroup-wide WGMMA instruction \cite[\S9.7.14]{ptx}—operates asynchronously and is capable of directly obtaining its input data from shared memory.

Hardware mechanisms enabling asynchrony make it possible to design warp-specialized kernels, wherein the warps of a given CTA are explicitly assigned to either producer or consumer tasks—issuing only data transfer or computation instructions, respectively. This approach broadly enhances the compiler's effectiveness at crafting efficient instruction schedules \citep{warp-specialization-2011}. Moreover, Hopper provides support for dynamically reallocating registers among warpgroups through the \verb|setmaxnreg| primitive \cite[\S9.7.17.1]{ptx}, allowing warps engaged in MMA operations to receive increased RMEM allocations relative to those limited to TMA tasks (which only require a single participating thread).

\paragraph{Low-precision number formats:}
\label{sec:low-precision-gpu}
Recent GPUs incorporate dedicated hardware to accelerate operations on low-precision data types. For instance, the WGMMA instruction exploits FP8 Tensor Cores on the Hopper architecture, providing double the per-SM throughput in comparison to computations performed in FP16 or BF16 precision.

Nonetheless, proper utilization of FP8 WGMMA requires careful attention to operand memory layout constraints. When performing a GEMM operation multiplying $A \times B^{\top}$, where $A$ is an $M\times K$ matrix and $B$ an $N\times K$ matrix, we classify an operand as \emph{mn-major} if its data is stored contiguously along the outer $M$ or $N$ axis, and as \emph{k-major} if contiguity is along the inner $K$ dimension. For FP16 WGMMA, both mn-major and k-major formats are permitted for input operands located in SMEM. In contrast, FP8 WGMMA restricts input formats to only k-major operands. This limitation becomes particularly challenging in scenarios such as attention mechanisms, where back-to-back GEMM operations are fused within a single kernel: mismatched FP32 accumulator layouts and FP8 operand layouts complicate the use of sequential, dependent FP8 WGMMA calls.

In the context of attention, these layout restrictions entail certain modifications to the design of an FP8 algorithm, which we describe in \cref{sec:algofp8}.

\subsection{Standard Attention and Flash Attention} As described by \citet{dao2022flashattention}, \textbf{standard attention} refers to the typical attention mechanism implemented on a GPU, where the intermediate matrices $\mathbf{S}$ and $\mathbf{P}$ are stored in High Bandwidth Memory (HBM). \textsc{FlashAttention}, on the other hand, is built around the concept of replacing these costly memory operations by introducing a localized softmax reduction, thereby streamlining attention computation into a unified kernel and minimizing intermediate data movement. In this context, the local softmax operation is executed in lines \ref{code-ws:softmax_start}-\ref{code-ws:softmax_end} of the consumer mainloop presented in \cref{alg:flash3_wgmma_ws_only}, along with blockwise rescaling of $\mathbf{O}$. For a straightforward derivation verifying that this method yields the correct computation of $\mathbf{O}$, see \cite[\S 2.3.1]{dao2023flashattention2}.

\section{FlashAttention-3: Algorithm}
\label{sec:algo}

This section presents the \textsc{FlashAttention-3} algorithm. Our exposition primarily addresses the forward pass, while the method for the backward pass is outlined in~\cref{sec:algo_ws_bwd}. We begin by detailing the integration of warp-specialization, utilizing a circular SMEM buffer, within the foundational structure of \textsc{FlashAttention-2}. Following this, we discuss leveraging WGMMA’s asynchrony to construct a two-stage pipeline where GEMM and softmax operations are overlapped. Lastly, we specify the required adaptations for FP8, covering both layout compliance and precision considerations through block quantization and incoherent processing techniques.

\subsection{Producer-Consumer asynchrony through warp-specialization and pingpong scheduling}
\label{sec:algo_ws}

\paragraph{Warp-specialization}
Similar to \textsc{FlashAttention-2}, the forward pass in \textsc{FlashAttention-3} achieves a high degree of parallelism across the batch size, number of heads, and the length of the query sequence. Therefore, we present a compute thread array (CTA)-oriented description of the method, which processes a tile $\mathbf{Q}_i$ from the query matrix and produces the corresponding tile $\mathbf{O}_i$ for the output. For clarity, we first introduce the warp-specialization approach assuming a circular SMEM buffer and excluding any GEMM-softmax overlap. Let $d$ denote the head dimension, $N$ the length of the sequence, and set a query block size $B_r$ such that $\mathbf{Q}$ is partitioned into $T_r = \lceil \frac{N}{B_r} \rceil$ blocks $\mathbf{Q}_1, .., \mathbf{Q}_{T_r}$.

\begin{algorithm}[H]
    \caption{\small\label{alg:flash3_wgmma_ws_only}\textsc{FlashAttention-3} forward pass \textbf{excluding} intra-consumer overlap -- CTA perspective}
    \begin{algorithmic}[1]
\REQUIRE Input matrices $\mathbf{Q}_i \in \mathbb{R}^{B_r \times d}$ along with $\mathbf{K}, \mathbf{V} \in \mathbb{R}^{N \times d}$ resident in HBM, key block size parameter $B_c$, and $T_c = \lceil \frac{N}{B_c} \rceil$.
\STATE Set up a pipeline construct for handling barrier synchronization via a circular SMEM buffer composed of $s$ stages.
\IF {executing within the producer warpgroup}
\STATE Free a fixed quantity of registers according to the allocation policy.
\STATE Transfer $\mathbf{Q}_i$ from HBM to shared memory.
\STATE Once the data is in shared memory, notify consumers that $\mathbf{Q}_i$ is ready.
\FOR{$0 \le j < T_c$}
    \STATE Await that the $(j\,\%\,s)$ buffer stage has been processed and released by consumers.
    \STATE Load $\mathbf{K}_j, \mathbf{V}_j$ from HBM to the $(j\,\%\,s)$ stage in shared memory.
    \STATE After completion, signal to consumers that $\mathbf{K}_j, \mathbf{V}_j$ have been loaded.
\ENDFOR
\ELSE
\STATE Allocate registers based on the expected count of consumer warps.
\STATE Initialize, on-chip, $\mathbf{O}_i = (0) \in \mathbb{R}^{B_r \times d}$; set $\ell_i, m_i = (0), (-\infty) \in \mathbb{R}^{B_r}$.
\STATE Wait until $\mathbf{Q}_i$ is available in shared memory.
\FOR{$0 \le j < T_c$}
\STATE Wait until $\mathbf{K}_j$ is present in shared memory.
\STATE Perform SS-GEMM: $\mathbf{S}_i^{(j)} = \mathbf{Q}_i \mathbf{K}_j^T$ and synchronize. \label{alg:ws_only_gemm1}
\STATE Store previous $m_i$ as $m_i^{\mathrm{old}}$, update $m_i$ to $\mathrm{max}(m_i^{\mathrm{old}}, \mathrm{rowmax}(\mathbf{S}_i^{(j)}))$. \label{code-ws:softmax_start}
\STATE Calculate normalized probabilities via $\widetilde{\mathbf{P}}_i^{(j)} = \mathrm{exp}(\mathbf{S}_i^{(j)} - m_i)$ and update $\ell_i$ with $\ell_i = \mathrm{exp}(m_i^{\mathrm{old}} - m_i) \ell_i + \mathrm{rowsum}(\widetilde{\mathbf{P}}_i^{(j)})$. \label{code-ws:softmax_end}
\STATE Ensure $\mathbf{V}_j$ has been loaded into shared memory.
\STATE Compute $\mathbf{O}_i = \mathrm{diag}(\mathrm{exp}(m_i^{\mathrm{old}} - m_i))^{-1} \mathbf{O}_i + \widetilde{\mathbf{P}}_i^{(j)} \mathbf{V}_j$ by RS-GEMM; synchronize. \label{alg:ws_only_gemm2}
\STATE Mark the $(j\,\%\,s)$ buffer stage as available to producers.
\ENDFOR
\STATE Finalize $\mathbf{O}_i = \mathrm{diag}(\ell_i)^{-1} \mathbf{O}_i$ and compute $L_i = m_i + \log(\ell_i)$.
\STATE Store $\mathbf{O}_i$ and $L_i$ in HBM as the $i$th portions of $\mathbf{O}$ and $L$.
\ENDIF
\end{algorithmic}
\end{algorithm}

In our realization of \cref{alg:flash3_wgmma_ws_only} targeting Hopper, we leverage \verb|setmaxnreg| for managing register (de)allocations, TMA to facilitate the loading of $\mathbf{Q}_i$ and $\{ \mathbf{K}_j, \mathbf{V}_j \}_{0 \leq j < T_c}$, and employ WGMMA to carry out the GEMM computations within the consumer mainloop. The prefixes SS or RS specify whether the initial operand resides in shared memory or within the register file, respectively. It is important to recognize, when following the execution path described by \cref{alg:flash3_wgmma_ws_only}, that TMA load instructions operate asynchronously and therefore do not block the progress of other load operations. Additionally, in the producer mainloop, wait operations are omitted during the initial $s$ cycles, allowing the buffer to be populated without interruption.

\paragraph{Pingpong scheduling} The asynchronous execution capabilities of WGMMA and TMA, combined with warp-specialized roles, make it possible to interleave the softmax computation for one warpgroup with the GEMM operations of another. This scheduling opportunity arises from a significant performance asymmetry: on contemporary hardware accelerators, matmul instructions greatly outpace the throughput of non-matmul functions. For instance, the H100 SXM5 GPU achieves a peak of 989 TFLOPS for FP16 matmul operations, yet special function throughput (such as exponentials required for softmax) is limited to just 3.9 TFLOPS\footnote{The CUDA programming guide indicates that each streaming multiprocessor (SM) is capable of executing 16 special function operations per clock cycle. Calculating 16 operations per SM, with 132 SMs running at 1830 MHz yields a total of 3.9 TFLOPS for special functions.}. In the FP16 attention forward path with a 128-dimensional head, the number of required matmul FLOPS outstrips exponentials by a factor of 512, while the hardware executes exponentials at 256 times lower throughput, causing the softmax computation to consume nearly half as many cycles as the matmul phase. The discrepancy is exacerbated with FP8: while matmul throughput doubles, the throughput for exponentials remains unchanged.

Because the exponential function is computed by a dedicated hardware component (the multi-function unit), optimal utilization entails initiating the exponential computation concurrently with the Tensor Cores' execution of matrix multiplications. To coordinate this, we leverage synchronization barriers (\texttt{bar.sync} instructions) to ensure that the GEMMs (specifically, GEMM1 -- $\mathbf{P} \mathbf{V}$ for a given iteration, and GEMM0 -- $\mathbf{Q} \mathbf{K}^\top$ for the following iteration) assigned to warpgroup 1 are executed prior to the corresponding GEMMs of warpgroup 2. This arrangement permits the softmax computation for warpgroup 1 to take place while warpgroup 2 processes its matrix multiplications. Subsequently, the warpgroups alternate roles: warpgroup 2 computes softmax while warpgroup 1 proceeds with its GEMMs, yielding a ``pingpong'' execution schedule as depicted in~\cref{fig:pingpong_scheduling}. Although real-world scheduling exhibits greater complexity than the idealized scenario illustrated, this approach consistently yields performance gains (for instance, throughput increases from 570 TFLOPS to 620-640 TFLOPS for FP16 forward passes with head size 128 and sequence length 8192).

\paragraph{Attention variants} For architectures such as multi-query attention~\citep{shazeer2019fast} and grouped query attention~\citep{ainslie2023gqa}, we adopt the methodology described in \textsc{FlashAttention-2}, modifying tensor indexing strategies to prevent redundant storage of $\mathbf{K}$ and $\mathbf{V}$ within HBM.

\subsection{Intra-warpgroup overlapping GEMMs and softmax}

Even within one warpgroup, we can overlap some instructions in the softmax with some instructions in the GEMMs. We describe one technique to do so.

Within the attention algorithm, the primary loop is constrained by intrinsic sequential dependencies, which limit opportunities for intra-iteration parallelism. Specifically, the computation of (local) softmax (lines \ref{code-ws:softmax_start} to \ref{code-ws:softmax_end}) is contingent upon the output $\mathbf{S}_i^{(j)}$ generated by the initial GEMM, and the subsequent GEMM depends on the intermediate value $\widetilde{\mathbf{P}}_i^{(j)}$. As a result, the execution order is enforced by wait statements at lines \ref{alg:ws_only_gemm1} and \ref{alg:ws_only_gemm2} in \cref{alg:flash3_wgmma_ws_only}, causing serialization among softmax and GEMM operations. To alleviate these restrictions, inter-iteration pipelining can be achieved by allocating supplemental register buffers, thereby breaking the serial dependencies. In this context, we introduce a GEMM-softmax pipelining method utilizing two stages\footnote{We remark that the number of overlapping pipeline stages may not necessarily match $s$, the number of stages defined by the circular SMEM buffer, although it does not exceed it.}:

\begin{algorithm}[H]
    \caption{\small\label{alg:flash3_wgmma}\textsc{FlashAttention-3} Consumer Warpgroup Forward Pass}
    \begin{algorithmic}[1]
      \REQUIRE  Inputs: Matrices $\mathbf{Q}_i \in \mathbb{R}^{B_r \times d}$ and $\mathbf{K}, \mathbf{V} \in \mathbb{R}^{N \times d}$ located in HBM; key block size $B_c$ with $T_c = \lceil \frac{N}{B_c} \rceil$.
      \STATE Dynamically reallocate a designated number of registers depending on how many consumer warps are active.
      \STATE Initialize on chip: set $\mathbf{O}_i = (0) \in \mathbb{R}^{B_r \times d}$, $\ell_i = 0$, and $m_i = -\infty$, all in $\mathbb{R}^{B_r}$.
      \STATE Ensure both $\mathbf{Q}_i$ and $\mathbf{K}_0$ are ready in shared memory before proceeding.
      \STATE Utilize WGMMA to calculate $\mathbf{S}_{\mathrm{cur}} = \mathbf{Q}_i \mathbf{K}_0^T$. After submitting, wait until completion.
      \STATE Deallocate the first buffer stage associated with $\mathbf{K}$.
      \STATE Determine $m_{i}$, $\tilde{\mathbf{P}}_{\mathrm{cur}}$, and $\ell_{i}$ from $\mathbf{S}_{\mathrm{cur}}$, and adjust $\mathbf{O}_i$ accordingly.
      \FOR{$1 \le j < T_c - 1$}
        \STATE Ensure that $\mathbf{K}_j$ is present in shared memory before proceeding.
        \label{code:mainloop_start}
        \STATE Invoke WGMMA for $\mathbf{S}_{\mathrm{next}} = \mathbf{Q}_i \mathbf{K}_{j}^T$, issue the computation, but continue without waiting. \label{code:first_wgmma}
        \STATE Await availability of $\mathbf{V}_{j-1}$ in shared memory.
        \STATE Update the running output with $\mathbf{O}_{i} = \mathbf{O}_{i} + \tilde{\mathbf{P}}_{\mathrm{cur}} \mathbf{V}_{j-1}$ via WGMMA, again committing the operation but not waiting. \label{code:second_wgmma}
        \STATE Wait until the computation for WGMMA $\mathbf{Q}_i \mathbf{K}_{j}^T$ is finished.
        \STATE Calculate $m_{i}$, $\tilde{\mathbf{P}}_{\mathrm{next}}$, and $\ell_{i}$ using results from $\mathbf{S}_{\mathrm{next}}$. \label{code:softmax}
        \STATE Following WGMMA completion for $\tilde{\mathbf{P}}_{\mathrm{cur}} \mathbf{V}_{j-1}$, update the scale of $\mathbf{O}_i$ as required.
        \STATE Free the $(j\,\%\,s)$th stage of the $\mathbf{K}$ buffer, as well as the $(j-1\,\%\,s)$th stage of the $\mathbf{V}$ buffer.
        \STATE Move $\mathbf{S}_{\mathrm{next}}$ to $\mathbf{S}_{\mathrm{cur}}$ to prepare for the next loop.
        \label{code:mainloop_end}
      \ENDFOR \STATE Wait until $\mathbf{V}_{T_c - 1}$ is present in shared memory.
      \STATE Compute the final update:
      $\mathbf{O}_{i} = \mathbf{O}_{i} + \tilde{\mathbf{P}}_{\mathrm{last}} \mathbf{V}_{T_c - 1}$ using WGMMA, waiting until operation is completed.

\STATE Epilogue: Adjust $\mathbf{O}_{i}$ by rescaling according to $m_{i}$. Derive $L_{i}$ utilizing both $m_{i}$ and $\ell_{i}$. Store the updated $\mathbf{O}_{i}$ and $L_{i}$ into HBM, populating the $i$-th sections of $\mathbf{O}$ and $L$ respectively.
    \end{algorithmic}
  \end{algorithm}

  \cref{alg:flash3_wgmma} acts as a substitute for the consumer sequence described in \cref{alg:flash3_wgmma_ws_only}, enabling the full implementation of the \textsc{FlashAttention-3} procedure in FP16 format. On a conceptual level, the term WGMMA serves as shorthand for asynchronous GEMM execution. During the loop spanning lines \ref{code:mainloop_start} through \ref{code:mainloop_end}, the second instance of the WGMMA call within each loop at step $j$ (line \ref{code:second_wgmma}) is executed in parallel with softmax computations pertaining to the subsequent iteration, $j+1$ (line \ref{code:softmax}).

Although the pipelined design outlined previously promises significant improvements in theoretical throughput, several pragmatic concerns arise in its implementation:
\paragraph{Compiler reordering} While the pseudocode is intended to enforce a specific pipeline between WGMMA and non-WGMMA operations, real-world compilation—especially by the NVCC compiler—may alter instruction sequencing during optimization phases. Such reordering could undermine the intended instruction overlap, causing either unanticipated effects or a reduction in the expected performance benefits. Examination of the resulting SASS code confirms that, in practice, the compiler does generate overlapped operations in line with the pipeline strategy (see Section~\ref{sec:2-stage-sass}).
\paragraph{Register pressure} Achieving high efficiency is contingent upon minimizing register spilling. The necessity of additional storage for intermediate pipeline values in the 2-stage architecture places further pressure on available registers. For instance, a separate $\mathbf{S}_{\mathrm{next}}$ value needs to be retained in registers, incurring an overhead of $B_r \times B_c \times \text{sizeof}(\text{float})$ registers per threadblock. This additional register use can interfere with other optimizations, such as employing larger block dimensions, which themselves place heavy demands on register resources. Practical implementation should therefore be guided by profiling to determine appropriate compromises.
\paragraph{3-stage pipelining} Building on the previously discussed 2-stage approach, a 3-stage pipeline is suggested to overlap a second WGMMA phase with softmax computation, which theoretically allows for even greater utilization of Tensor Cores. Nonetheless, this comes at the cost of even higher register usage owing to the further increased pipeline depth. Consequently, managing the balance between maximizing tile size and increasing pipeline stages becomes increasingly complex. A comprehensive explanation of the 3-stage strategy and related empirical findings is provided in ~\cref{sec:3-stage}.

\subsection{Low-precision with FP8}
\label{sec:algofp8}

\textbf{Efficiency: layout transformations.} Executing the forward pass of \textsc{FlashAttention-3} in FP8 precision introduces additional difficulties in layout handling that are not present when using FP16.

To elaborate, input tensors $\mathbf{Q}$, $\mathbf{K}$, and $\mathbf{V}$
are usually stored such that their data is contiguous along the head dimension. However, in order to meet the k-major layout requirement imposed by the FP8 WGMMA kernel during the second GEMM, it is necessary for $\mathbf{V}$—specifically, for the $\mathbf{V}$ tiles loaded into SMEM—to be contiguous along the sequence length dimension instead. Since TMA loads are unable to alter which dimension is contiguous, there are two available strategies: (1) perform a full transpose of $\mathbf{V}$ within GMEM prior to the GEMM, or (2) apply a tile-wise transpose to portions of $\mathbf{V}$ inside the kernel, following their loading into SMEM. For the first approach, one may either (1a) incorporate the transpose operation directly into the epilogue of a prior step (such as the rotary embedding) in a fused fashion, or (1b) launch a dedicated pre-processing transpose kernel\footnote{An efficient transpose kernel can reach performance close to the device's memory bandwidth~\citep{colfax_cutlass_transpose_2024}.}, swapping the head and sequence length strides. Nevertheless, (1a) is impractical to generalize as part of a typical library, and (1b) is prohibitively inefficient in inference, which is typically limited by memory bandwidth.

For FP8 \textsc{FlashAttention-3}, we select strategy (2). During the in-kernel transposition process, we utilize the LDSM (\verb|ldmatrix|) and STSM (\verb|stmatrix|) instructions, which enable a warp of threads to collaboratively transfer data between shared memory (SMEM) and registers (RMEM) in blocks of 128 bytes.\footnote{While the PTX documentation specifies that LDSM/STSM handle $8 \times 8$ matrices composed of 16-bit elements \cite[\S 9.7.13.4.15-16]{ptx}, it is possible to pack pairs of 8-bit values and apply LDSM/STSM for computations in FP8. Nevertheless, the transposing modes of LDSM/STSM lack the capability to separate these packed 8-bit elements, thereby requiring additional register rearrangements between LDSM and STSM operations to achieve tile-level transposition; further elaboration is omitted here.} These instructions make efficient use of register resources, permitting execution directly in the producer warpgroup, and can alter memory layouts during data transfer. Furthermore, after the initial iteration, it is feasible to overlap the transposition of the upcoming $\mathbf{V}$ tile with the execution of the two WGMMAs that operate on the previous $\mathbf{V}$ and the current $\mathbf{K}$ tile.

Second, it is notable that, in contrast to FP16, the arrangement of the FP32 accumulator in FP8 WGMMA does not match the expected register layout of operand A. We illustrate segments of both memory organizations in \cref{fig:rmem_accum} and \cref{fig:rmem_operand}, showing how values are assigned to registers per thread in their respective orders. Employing byte permute instructions, we are able to reformat the output accumulator from the first WGMMA to conform to the requirements of the second WGMMA, aligning as well with the memory layout used for the $\mathbf{V}$ tile by the in-kernel transpose. In particular, referring to \cref{fig:rmem_accum}, we reorganize the ordering sequentially as
$$\{ \verb|d0 d1 d4 d5 d2 d3 d6 d7| \},$$
and replicate this register order for every group of 8 bytes. This effectively reorders the columns within the logical representation of the $\mathbf{P}$ tile (for instance, columns $0189$ are shifted to the lead columns). To guarantee that WGMMA generates the intended output tile, we can adjust the in-kernel transpose to output the $\mathbf{V}$ tile with a corresponding permutation applied to its rows.\footnote{This flexibility introduced by implementing the in-kernel transpose obviates the necessity of shuffle instructions for redistributing register contents among threads, as was previously required in~\citep{colfax_fp8_flashattention_2024}.}

\textbf{Accuracy: block quantization and incoherent processing.} The FP8 (e4m3) representation allocates 3 bits to the mantissa and 4 bits to the exponent, which leads to greater quantization error relative to FP16 or BF16 formats. Additionally, large-scale neural models often exhibit extreme outlier values~\citep{dettmers2208llm,
  sun2024massive} that are orders of magnitude larger than the majority of elements, posing challenges for effective quantization. Common practice involves per-tensor scaling~\citep{micikevicius2022fp8}, maintaining a single scale parameter for each tensor (such as $\mathbf{Q}$, $\mathbf{K}$, or $\mathbf{V}$).
To mitigate the accumulation of numerical errors in FP8 attention calculations, we adopt two approaches:

\begin{enumerate}

\item \textbf{Block quantization}: In this approach, a single scalar is retained for each block, meaning that for the tensors $\mathbf{Q}$, $\mathbf{K}$, and $\mathbf{V}$, we partition them into blocks with dimensions $B_r \times d$ or $B_c \times d$ and perform quantization separately on each block. Notably, this quantization step can be integrated with operations such as rotary embedding that precede the attention mechanism without introducing computational latency, as such operations are typically constrained by memory bandwidth. Since the \textsc{FlashAttention-3} algorithm inherently processes data in blocks, it is straightforward to rescale each block of $\mathbf{S}$ to incorporate block-level quantization, incurring no additional computational overhead.

\item \textbf{Incoherent processing}: To mitigate the impact of outliers, we apply a random orthogonal matrix $\mathbf{M}$ to both $\mathbf{Q}$ and $\mathbf{K}$ before quantizing to FP8. Given the orthogonality property $\mathbf{M} \mathbf{M}^\top = I$, the transformation preserves the attention computation, i.e., $(\mathbf{Q} \mathbf{M})(\mathbf{K} \mathbf{M})^\top = \mathbf{Q} \mathbf{K}^\top$. This procedure redistributes outlier values, as each component of $\mathbf{Q}\mathbf{M}$ or $\mathbf{K}\mathbf{M}$ becomes a random linear combination of the elements in $\mathbf{Q}$ or $\mathbf{K}$, ultimately lowering the quantization error. Consistent with the methodology in \citet{chee2024quip} and~\citet{tseng2024quip}, $\mathbf{M}$ is selected as the product of random diagonal matrices with entries $\pm1$ and a Hadamard matrix, enabling efficient computation in $O(d \log d)$ time rather than $O(d^2)$, and this transformation can be combined with rotary embedding seamlessly, incurring no further computational cost.
\end{enumerate}

We validate that these two techniques reduces numerical error by up to 2.6$\times$ in \cref{sec:numerical_error}.

\section{Empirical Validation}
\label{sec:exp}

We use the primitives from CUTLASS~\citep{Thakkar_CUTLASS_2023} such as WGMMA and TMA abstractions to implement \textsc{FlashAttention-3} and evaluate its efficiency and accuracy.

\begin{itemize}

\item \textbf{Benchmarking attention.} We evaluate the execution time of \textsc{FlashAttention-3} at various sequence lengths, contrasting its performance with the baseline PyTorch implementation, \textsc{FlashAttention-2}, \textsc{FlashAttention-2} implemented in Triton (leveraging H100-specific capabilities), as well as a cuDNN-based vendor implementation of \textsc{FlashAttention-2} optimized for H100 GPUs. Our results indicate that \textsc{FlashAttention-3} achieves speedups of up to 2.0$\times$ compared to \textsc{FlashAttention-2} and 1.5$\times$ compared to its Triton variant. Furthermore, \textsc{FlashAttention-3} attains a throughput of up to 740 TFLOPs/s, amounting to 75\% of the H100 GPU’s peak theoretical TFLOPs/s.
\item \textbf{Ablation study.} Through controlled experiments, we demonstrate that enhancements such as warp-specialization and GEMM-softmax pipelining significantly accelerate \textsc{FlashAttention-3}.
\item \textbf{Accuracy of FP8 attention.} Our analyses show that techniques including block quantization and incoherent processing reduce the numerical error in FP8 \textsc{FlashAttention-3} by a factor of 2.6$\times$.
\end{itemize}

\subsection{Benchmarking Attention}
\label{subsec:benchmark_attn}

We evaluate the execution times of several attention mechanisms on an H100 80GB SXM5 GPU under various configurations—both with and without a causal mask, and head dimensions set to 64 or 128—using FP16 precision for the input data. The experimental findings are depicted in~\cref{fig:benchmark_attn_fwd} and~\cref{fig:benchmark_attn_bwd}. Our results indicate that \textsc{FlashAttention-3} achieves speedups of approximately 1.5-2.0$\times$ over \textsc{FlashAttention-2} during the forward pass and about 1.5-1.75$\times$ in the backward pass. When benchmarked against a standard attention implementation, \textsc{FlashAttention-3} reaches acceleration factors ranging from 3$\times$ up to 16$\times$. Notably, for sequence lengths of 1,000 and greater, \textsc{FlashAttention-3} even outperforms a proprietary vendor library (cuDNN—closed source), which is specifically optimized for H100 GPUs.

\paragraph{Benchmark settings:} The sequence length is adjusted through values of 512, 1k, up to 16k, while the batch size is chosen such that the cumulative token count remains 16k. The hidden dimension is fixed at 2048, and the head dimension is tested with values of 64, 128, and 256, corresponding to configurations with 32, 16, and 8 attention heads respectively. The FLOPs for the forward pass are computed as follows:

\begin{equation*}
  4 \cdot \text{seqlen}^2 \cdot \text{head dimension} \cdot \text{number of heads}.
\end{equation*}

Applying causal masking reduces the total count by half, since roughly 50\% of the computations are performed. For the backward pass, the FLOPs are obtained by scaling the forward pass FLOPs by a factor of 2.5 (this is because the backward computation involves 5 matrix multiplications compared to the 2 performed during the forward pass, owing to the necessity of recomputation).

Additionally, we evaluate the forward pass runtime for FP8 precision using similar experimental conditions. Outcomes for headdim 256 are presented in ~\cref{fig:benchmark_attn_fp8}, with a comprehensive breakdown available in \cref{sec:benchmark_attn_fp8_full}.

\subsection{Ablation Study: 2-Stage Pipelining Experiments}
\label{sec:2-stage-pipelining-experiments}

We conduct an ablation study on both the 2-stage WGMMA-softmax pipelining and warp-specialization for non-causal FP16 \textsc{FlashAttention-3}, utilizing fixed settings of $\{\text{batch}, \text{seqlen}, \text{nheads}, \text{hdim}\} = \{ 4, 8448, 16, 128\}$. The outcomes presented in~\cref{table:ablation_pipelining} demonstrate that incorporating our algorithmic innovations—namely, asynchronous execution through warp-specialization and the concurrent processing of GEMM and softmax—achieves notable performance gains, improving throughput from 570 to 661 TFLOPs.

\subsection{Numerical Error Validation}
\label{sec:numerical_error}

Given the ongoing attention to the numerical errors present in \textsc{FlashAttention}~\citep{golden2024flash}, we evaluate and contrast \textsc{FlashAttention-2}, \textsc{FlashAttention-3}, and a conventional attention implementation by benchmarking them against a reference in FP64 precision. In order to mimic the outlier activations and features observed in LLMs~\citep{dettmers2208llm, sun2024massive}, we sample the elements of $\mathbf{Q}, \mathbf{K}, \mathbf{V}$ according to the following distribution:

\begin{equation*}
  \mathcal{N}(0, 1) + \mathcal{N}(0, 100) \cdot \mathrm{Bernoulli}(0.001).
\end{equation*}

In other words, all entries follow a normal distribution with a mean of zero and standard deviation of 1, except for 0.1\% of the entries, for which an additional independent normally distributed variable with standard deviation 10 is included. We calculate the root mean squared error (RMSE) as shown in~\cref{table:numerical_error}. Using FP16, both \textsc{FlashAttention-2} and \textsc{FlashAttention-3} yield an RMSE that is 1.7$\times$ lower than that of the standard approach, due to retaining intermediate computations (specifically, the softmax outputs) in FP32 precision. In the FP8 baseline, per-tensor scaling is used, the matrix multiplication accumulator operates in FP32, and the softmax intermediates are stored in FP16. Benefiting from block quantization and incoherent computation, \textsc{FlashAttention-3} in FP8 attains a 2.6$\times$ improvement in accuracy over this baseline.

\section{Dicussion, Limitations, Conclusion}
\label{sec:discussion}

Through \textsc{FlashAttention-3}, we show that leveraging novel programming paradigms and harnessing hardware capabilities like asynchrony and low-precision computation can significantly enhance both the performance and precision of attention mechanisms. Our approach yields a 1.5-2.0$\times$ improvement in attention speed over \textsc{FlashAttention-2} and achieves a 2.6$\times$ reduction in FP8 numerical errors when compared to typical per-tensor quantization. Future work aims to address certain constraints: further optimization for LLM inference, integrating persistent kernel methodologies within the FP8 kernel,\footnote{In our evaluations, the FP16 variant of \textsc{FlashAttention-3} incorporates a persistent kernel and implements load balancing, while this is absent in the FP8 variant. This partially accounts for the reduced performance of FP8 \textsc{FlashAttention-3} on tasks with shorter sequence lengths and causal masking, relative to FP8 cuDNN kernels.} and investigating the implications of low-precision attention during large-scale model training. Although our present research centers on Hopper GPUs, we anticipate that the strategies introduced here can generalize to a wide array of hardware accelerators. Ultimately, we envision that more efficient and precise attention primitives will pave the way for advances in long-context application domains.

\end{document}